\Needspace{4\baselineskip}
\begin{Example}
\label{ex:abel-integral-equation-laplace}
Abel 积分方程为
\begin{equation}\label{eq:abel-integral-equation}
\int_0^x\dfrac{y(t)}{\sqrt{x-t}}\,\mathrm{d}t=\sqrt\pi f(x),
\end{equation}
仍然寻求 $x>0$ 时的解, 并认为 $y$ 和 $f$ 对负的 $x$ 作了零延拓. 为使解存在, 显然必须有 $f(0)=0$. 左端为 $y$ 与 $x_+^{-1/2}$ 的卷积, 而卷积的 Laplace 变换是各 Laplace 变换的乘积. 为求 $x_+^{-1/2}$ 的变换, 计算
\begin{equation}\label{eq:laplace-inverse-square-root-kernel}
\int_0^\infty e^{-sx}x^{-1/2}\,\mathrm{d}x=\dfrac1{\sqrt s}\int_0^\infty e^{-t}t^{-1/2}\,\mathrm{d}t=\sqrt{\dfrac\pi s}
\end{equation}
对任意实的正数 $s$ 成立, 并由解析延拓的唯一性, 对复数 $s$ 也成立. 因而变换后的方程为
\begin{equation}\label{eq:abel-equation-laplace-transform}
\dfrac{\mathcal{L}[y](s)}{\sqrt s}=\mathcal{L}[f](s),
\end{equation}
将它写为
\begin{equation}\label{eq:abel-equation-transform-solution}
\mathcal{L}[y](s)=\dfrac{s\mathcal{L}[f](s)}{\sqrt s}.
\end{equation}
作逆变换, 得到
\begin{equation}\label{eq:abel-equation-inverse-transform-solution}
y(x)=\dfrac1{\sqrt\pi}\int_0^x\dfrac{f'(t)}{\sqrt{x-t}}\,\mathrm{d}t.
\end{equation}

\end{Example}

\Needspace{4\baselineskip}
\begin{Example}
\label{ex:heat-boundary-laplace}
Laplace 变换也适用于偏微分方程的初值问题. 首先指出, Definition~\ref{def:tempered-fourier-transform} 容易推广, 以定义广义函数仅关于一部分变量的 Fourier 变换. 例如, 处理初值问题时, 可以关于时间作 Laplace 变换. 当然, 要使边界条件有意义, 对解的了解必须多于仅仅知道它是广义函数. 例如, 在以下问题中, 可以将 $u$ 视为以 $x$ 为参数的 $t$ 的广义函数.

考虑初边值问题
\begin{equation}\label{eq:heat-unit-boundary-zero-initial-problem}
\begin{aligned}
u_t&=u_{xx},\quad x\in(0,1),\ t>0,\\
u(x,0)&=0,\quad x\in(0,1),\\
u(0,t)&=u(1,t)=1,\quad t>0.
\end{aligned}
\end{equation}
如往常一样, 对负的 $t$ 将 $u$ 作零延拓. 关于时间的 Laplace 变换给出问题
\begin{equation}\label{eq:heat-unit-boundary-laplace-problem}
\begin{aligned}
s\mathcal{L}[u](x,s)&=\mathcal{L}[u]_{xx}(x,s),\\
\mathcal{L}[u](0,s)&=\mathcal{L}[u](1,s)=\dfrac1s.
\end{aligned}
\end{equation}
这一方程具有解
\begin{equation}\label{eq:heat-unit-boundary-laplace-solution}
\mathcal{L}[u](x,s)=\dfrac{\cosh(\sqrt s(x-1/2))}{s\cosh(\sqrt s/2)}.
\end{equation}
逆变换公式给出
\begin{equation}\label{eq:heat-unit-boundary-inverse-laplace-formula}
u(x,t)=\dfrac1{2\pi i}\int_{\mu-i\infty}^{\mu+i\infty}e^{st}\dfrac{\cosh(\sqrt s(x-1/2))}{s\cosh(\sqrt s/2)}\,\mathrm{d}s.
\end{equation}
这里 $\mu$ 可取任意正数. 该积分不能用闭式求出, 但当然可以数值计算; 也可以用它推断解的定性性质, 如正则性或当 $t\to\infty$ 时的渐近行为.
\end{Example}
